\documentclass[border=12pt]{standalone}
\usepackage[utf8]{inputenc}
\usepackage[T1]{fontenc}
\usepackage[american, siunitx]{circuitikz}

\sisetup{output-decimal-marker={,}}

\ctikzset{
    resistors/scale=0.8,
    diodes/scale=0.8,
    switches/scale=0.8
}

\begin{document}
\begin{circuitikz}[font=\sffamily\small]

    % =========================================================================
    % 1. BLOCO ESQUERDO: TECLADO ANALOGICO (ESCADA RESISTIVA EM A0)
    % =========================================================================
    % Alimentacao +5V no topo da escada resistiva
    \draw (-4.0, 4.4) node[vcc]{+5V}
        to[R, l^=R2, a_=2<\kilo\ohm>] (-4.0, 2.7)
        to[R, l^=R3, a_=330<\ohm>] (-4.0, 1.1)
        to[R, l^=R4, a_=620<\ohm>] (-4.0, -0.5)
        to[R, l^=R5, a_=1<\kilo\ohm>] (-4.0, -2.1)
        to[R, l^=R6, a_=3.3<\kilo\ohm>] (-4.0, -3.7);

    % 5 Botoes (S1 a S5) ligados entre cada no da escada e o barramento GND
    \foreach \y/\swname/\func in {
        2.7/S1/RIGHT,
        1.1/S2/UP,
        -0.5/S3/DOWN,
        -2.1/S4/LEFT,
        -3.7/S5/SELECT%
    }{
        \draw (-7.8, \y) to[push button, l^={\swname\ (\func)}] (-4.0, \y);
    }

    % Barramento comum de GND dos botoes (sem ponto no canto superior)
    \draw (-7.8, 2.7) -- (-7.8, -4.6) node[ground]{};

    % Pontos de juncao apenas nas conexoes em T (exclui cantos de 90 graus)
    \foreach \y in {2.7, 1.1, -0.5, -2.1}{
        \node[circ] at (-4.0, \y) {};
    }
    \foreach \y in {1.1, -0.5, -2.1, -3.7}{
        \node[circ] at (-7.8, \y) {};
    }

    % =========================================================================
    % 2. BLOCO CENTRAL: ARDUINO UNO R3 (U1)
    % =========================================================================
    \draw[thick, fill=gray!5] (0, -5.2) rectangle (4.0, 4.2);
    \node[align=center, font=\sffamily\bfseries] at (2.0, 0.3) {U1\\Arduino\\Uno R3};
    \node[font=\sffamily\scriptsize, text=gray!70!black] at (2.0, -0.7) {(ATmega328P)};

    % Alimentacao do Arduino
    \draw (2.0, 4.2) -- (2.0, 5.0) node[vcc]{+5V};
    \node[below, font=\ttfamily\footnotesize] at (2.0, 4.2) {5V};

    \draw (2.0, -5.2) -- (2.0, -6.0) node[ground]{};
    \node[above, font=\ttfamily\footnotesize] at (2.0, -5.2) {GND};

    % Entrada analogica A0 vinda da escada resistiva
    \draw (-4.0, 2.7) -- (0, 2.7);
    \node[right, font=\ttfamily\footnotesize] at (0, 2.7) {A0};

    % Saida PWM no pino D3
    \node[left, font=\ttfamily\footnotesize] at (4.0, -3.1) {D3};
    \draw[->, >=stealth, thick] (4.0, -3.1) -- (6.0, -3.1)
        node[right, font=\sffamily\bfseries\footnotesize] {PWM};

    % Saida LED no pino D13 (LED D1 + Resistor R8 de 220 Ohms para GND)
    \node[left, font=\ttfamily\footnotesize] at (4.0, -4.4) {D13};
    \draw (4.0, -4.4) -- (4.5, -4.4)
        to[leDo, l^=D1] (6.3, -4.4)
        to[R, l^=R8, a_=220<\ohm>] (8.2, -4.4)
        -- (8.2, -5.2) node[ground]{};

    % =========================================================================
    % 3. BLOCO DIREITO: DISPLAY LCD 16x2 (HD44780)
    % =========================================================================
    \draw[thick, fill=gray!5] (7.6, -2.8) rectangle (12.6, 3.8);
    \node[align=center, font=\sffamily\bfseries] at (10.3, 0.7) {LCD1\\Display LCD\\16$\times$2};
    \node[font=\sffamily\scriptsize, text=gray!70!black] at (10.3, -0.6) {(HD44780)};

    % Conexoes Arduino -> LCD (Modo 4 bits)
    \foreach \y/\apin/\lpin in {
        3.0/D8/RS (4),
        2.0/D9/E (6),
        1.0/D4/DB4 (11),
        0.0/D5/DB5 (12),
        -1.0/D6/DB6 (13),
        -2.0/D7/DB7 (14)
    }{
        \node[left, font=\ttfamily\footnotesize] at (4.0, \y) {\apin};
        \draw (4.0, \y) -- (7.6, \y);
        \node[right, font=\ttfamily\scriptsize] at (7.6, \y) {\lpin};
    }

    % Alimentacao superior do LCD (VDD e Backlight BL+ com resistor de 220R)
    \draw (9.2, 3.8) -- (9.2, 4.6) node[vcc]{+5V};
    \node[below, font=\ttfamily\scriptsize] at (9.2, 3.8) {VDD (2)};

    \draw (11.4, 3.8) to[R, l_=R7, a^=220<\ohm>] (11.4, 6.0) node[vcc]{+5V};
    \node[below, font=\ttfamily\scriptsize] at (11.4, 3.8) {BL+ (15)};

    % Aterramento inferior do LCD (VSS, RW e BL-)
    \node[above, font=\ttfamily\scriptsize] at (8.8, -2.8) {VSS (1)};
    \node[above, font=\ttfamily\scriptsize] at (10.2, -2.8) {RW (5)};
    \node[above, font=\ttfamily\scriptsize] at (11.6, -2.8) {BL- (16)};

    \draw (8.8, -2.8) -- (8.8, -3.8) -- (11.6, -3.8) -- (11.6, -2.8);
    \draw (10.2, -2.8) -- (10.2, -4.6) node[ground]{};
    \node[circ] at (10.2, -3.8) {};

    % Ajuste de Contraste (V0) com potenciometro RP1 (20 kOhm) espelhado
    \node[left, font=\ttfamily\scriptsize] at (12.6, 0.5) {V0 (3)};
    \draw (14.4, 2.2) node[vcc]{+5V}
        to[pR, mirror, n=pot, l^={\shortstack[l]{RP1\\\SI{20}{\kilo\ohm}}}] (14.4, -1.2) node[ground]{};
    \draw (12.6, 0.5) -- (pot.wiper);

\end{circuitikz}
\end{document}